\documentclass[conference,a4paper]{IEEEtran}
\IEEEoverridecommandlockouts
% The preceding line is only needed to identify funding in the first footnote. If that is unneeded, please comment it out.
%Template version as of 6/27/2024

\usepackage{cite}
\usepackage{amsmath,amssymb,amsfonts}
\usepackage{algorithmic}
\usepackage{graphicx}
\usepackage{textcomp}
\usepackage{xcolor}
\usepackage{booktabs}
\usepackage{url}
\def\BibTeX{{\rm B\kern-.05em{\sc i\kern-.025em b}\kern-.08em
    T\kern-.1667em\lower.7ex\hbox{E}\kern-.125emX}}
\begin{document}

\title{Assignment 1 : Supervised Machine Learning %\\
%{\footnotesize \textsuperscript{*}Note: Sub-titles are not captured for https://ieeexplore.ieee.org  and
%should not be used}
%\thanks{Identify applicable funding agency here. If none, delete this.}
}

\author{\IEEEauthorblockN{Renjiro Owen}
\textit{Victoria University Wellington}\\
Wellington, New Zealand\\
owenrenj@myvuw.ac.nz}

\maketitle

\section{Decision Tree Method}

\subsection{Results}

The decision tree algorithm was trained using the rtg A, rtg B, and rtg C training dataset. 
The generated decision tree was constructed using entropy-based Information Gain, 
with a stopping threshold of $IG < 0.00001$.
The overall training classification accuracy achieved was:

\begin{equation}
	Accuracy_{rtgA} = \text{100}\%
\end{equation}

\begin{equation}
	Accuracy_{rtgB} = \text{100}\%
\end{equation}

\begin{equation}
	Accuracy_{rtgC} = \text{100}\%
\end{equation}

The classification accuracy was calculated using:

\begin{equation}
	Accuracy = \frac{\text{Number of correctly classified instances}}
	{\text{Total number of instances}}
\end{equation}

\subsection{Minimal Entropy}

For each dataset, the feature resulting in the minimum entropy was identified by
calculating the weighted entropy after splitting the entire training dataset on
each available feature. The feature with the lowest resulting entropy represents
the split that produces the purest separation of class labels and therefore
provides the greatest reduction in uncertainty.

For \textit{rtg A}, Feature 0 produced the minimum entropy value of 0.285506 with
an Information Gain of 0.116673. Therefore, Feature 0 was selected as the root
splitting feature of the generated decision tree. For \textit{rtg B}, Feature 3
resulted in the lowest entropy value of 0.900495 with an Information Gain of
0.053199, making it the most effective feature for the initial split. For
\textit{rtg C}, Feature 3 achieved an entropy value of 0.000000 and an Information
Gain of 0.904381, indicating that this feature completely separates the classes
in the training dataset.

These observations were made by comparing the entropy values obtained from all
features. The feature with the smallest entropy (or equivalently the highest
Information Gain) was considered the most informative feature and was chosen by
the decision tree algorithm for splitting the dataset.

\subsection{Unused Feature}
The features utilized by the decision tree were different for each dataset. For
\textit{rtg A}, Feature 3 was not used in constructing the decision tree. This
occurred because the other features provided sufficient Information Gain to
separate the classes, allowing the tree to reach pure leaf nodes before Feature
3 was considered. Since the stopping criteria were satisfied once the nodes
became sufficiently pure, Feature 3 did not contribute additional useful
information for classification.

For \textit{rtg B}, all available features were utilized in the construction of
the decision tree. This indicates that each feature contributed useful
information during the splitting process. The tree required all features to
further separate the data until the stopping criteria were reached.

For \textit{rtg C}, only Feature 3 was used in constructing the decision tree,
while the remaining features were excluded. This was because Feature 3 provided
a very strong separation of the classes, achieving an entropy of 0 and a high
Information Gain of 0.904381. After splitting on Feature 3, the resulting nodes
were already pure, meaning additional features were unnecessary. Therefore, the
remaining features were not selected because they did not provide any further
improvement to the classification.

Overall, features were excluded from the tree-building process when they did not
provide enough additional Information Gain or when previous splits had already
produced pure leaf nodes. The decision tree only selects features that reduce
class uncertainty and ignores redundant or unnecessary attributes.

\subsection{Rtg C Discussion}
\subsubsection{Generalisation Ability of the Generated Tree}

The generated decision tree is not expected to perform well on unseen data. The
tree structure indicates that Feature 3 alone is able to completely separate the
training instances, resulting in a very simple tree where the split is entirely
dependent on a single feature. Although this produces perfect classification
performance on the training dataset, it may not generalise well to unseen data
because the model relies heavily on Feature 3.

If new instances contain feature values that were not represented in the training
dataset, the tree may not have sufficient information from other features to make
a reliable prediction. Therefore, the tree structure suggests a risk of poor
generalisation due to its dependency on a single feature rather than learning
patterns from multiple attributes.


\subsubsection{Unused Features in Tree Construction}

For the rtg C dataset, only Feature 3 was utilized in constructing the decision
tree, while the remaining features were not selected. This occurred because
Feature 3 provided a very strong separation between the two classes, achieving an
entropy of 0 and a high Information Gain of 0.904381.

Since splitting on Feature 3 resulted in pure leaf nodes, there was no need for
the decision tree to consider additional features. The other attributes may
contain useful information; however, their contribution was unnecessary after
Feature 3 had already completely separated the training data. Therefore, these
features were excluded because they did not provide additional Information Gain
during the tree-building process.


\subsubsection{Potential Preprocessing Improvements}

Although a test dataset is not available, several preprocessing approaches could
be considered to improve the generalisation ability of the generated tree.
Firstly, Feature 3 could be investigated further because the tree's complete
dependency on this single attribute may indicate that the dataset contains a
dominant or potentially artificial feature. Removing Feature 3 before training
could encourage the model to learn from other attributes and produce a more
generalised decision boundary.

Additionally, feature selection techniques could be applied to identify and
retain informative attributes while reducing dependency on a single feature.
Tree pruning or increasing the Information Gain stopping threshold could also be
considered to prevent the tree from creating overly specific splits based only on
the training data.

These preprocessing steps may reduce overfitting and improve the ability of the
decision tree to classify unseen instances.

\bibliographystyle{IEEEtran}   % or plain, unsrt, apalike, etc.
\bibliography{references}

\end{document}
